% Unidad U013. Biografía estructural completa del carácter exponencial.

\label{gfp:b:chap:biografia-completa-e}

\section{Sélecteur orbital fini : existence et unicité de l’émission compatible}

Soit \(\mathcal W=\mathbb F_3^6\). Pour une émission
\(U=(u_1,\ldots,u_6)\), le mot visible est
\[
 \Psi(U)=((-u_1)\bmod3,\ldots,(-u_6)\bmod3).
\]
Sur \(\mathcal W\), l’action diédrale est donnée par
\[
\begin{aligned}
 r(u_1,u_2,u_3,u_4,u_5,u_6)
 &=(u_3,u_1,u_2,u_5,u_6,u_4),\\
 s(u_1,u_2,u_3,u_4,u_5,u_6)
 &=(u_4,u_5,u_6,u_1,u_2,u_3).
\end{aligned}
\]
Pour \(w\in\mathcal W\), nous définissons
\[
 n(w)=\#\Psi^{-1}(w),
 \qquad
 M(w)=\sum_{U\in\Psi^{-1}(w)}\operatorname{mult}(U),
\]
et
\[
 \nu(w)=
 \bigl(
 \#\{j:w_j=0\},
 \#\{j:w_j=1\},
 \#\{j:w_j=2\}
 \bigr).
\]

\begin{theorem}[Sélection structurelle du canal exponentiel]
\label{gfp:b:thm:seleccion-estructural-e}
Parmi les \(43\) orbites du catalogue visible, il existe une unique orbite dont les
fibres satisfont simultanément :
\[
 n(w)=1,\qquad M(w)=144,\qquad
 \mathfrak D(w)=\{0,3,6\},\qquad
 \nu(w)=(2,3,1).
\]
Le calibre
\[
 \operatorname{diag}_c=6,
 \qquad \operatorname{card}=ES
\]
y sélectionne
\[
 \boxed{w_6^e=201101.}
\]
\end{theorem}

\begin{proof}
Les quatre conditions sont évaluées sur les \(243\) mots du catalogue,
non sur un échantillon. Une unique orbite satisfait le prédicat conjoint. Dans
cette orbite, une seule émission possède à la fois la diagonale \(6\) et l’orientation
\(ES\). L’énumération utilise uniquement des champs du catalogue APP--TRIT et ne
consulte pas de développement de \(e\).
\end{proof}

\section{Relèvement par blocs et compatibilité des frontières successives}

La récurrence \(L_0,L_0,L_1,L_1\) produit
\[
 w_{30}^{e}
 =201101|121221|102011|012222|102011,
\]
et la frontière conjointe apporte
\[
 u_5^e=021222.
\]
Le registre de trente-six trits devient
\[
 w_{36}^{e}
 =201101|121221|102011|012222|102011|021222.
\]

\section{Signature positionnelle décimale et carré de compatibilité}

La lecture initiale de douze chiffres est
\[
 \mathbf d_e^{(12)}
 =718281828459
 =7\,\underbrace{1828\,1828}_{\text{carré local}}\,459.
\]
L’égalité centrale est la concaténation
\[
 1828\mathbin{\|}1828,
\]
non un produit arithmétique ni le début supposé d’une période.

\begin{definition}[Carré positionnel]
Un mot \(d_1\cdots d_{12}\) contient un carré de longueur \(\ell\)
à la position \(p\) lorsque
\[
 d_p\cdots d_{p+\ell-1}
 =d_{p+\ell}\cdots d_{p+2\ell-1}.
\]
La position et les contextes latéraux appartiennent à la signature.
\end{definition}

\begin{table}[htbp]
\centering
\caption{Recensement exhaustif des carrés dans les lectures de douze chiffres.}
\label{gfp:b:tab:censo-cuadrados-e}
\begin{tabular}{c*{6}{r}}
\toprule
Longueur \(\ell\)&1&2&3&4&5&6\\
\midrule
Occurrences&240&21&3&2&0&0\\
Mots qui les contiennent&153&19&3&2&0&0\\
\bottomrule
\end{tabular}
\end{table}

L’autre mot contenant un carré de longueur quatre est
\[
 575157518472
 =\underbrace{5751\,5751}_{\text{position initiale }1}\,8472.
\]

\begin{proposition}[Unicité positionnelle]
\label{gfp:b:prop:cuadrado-posicional-e}
La lecture du canal \(e\) est la seule qui possède un carré de longueur
quatre après un préfixe de longueur un et avec les contextes \(7\) et \(459\).
\end{proposition}

\begin{proof}
Le recensement ne laisse que deux mots contenant des carrés de longueur quatre. L’un
commence son carré à la première position ; l’autre après le préfixe \(7\). La
position et les contextes séparent les deux possibilités.
\end{proof}

\section{Retour de phase avec changement d’état et conservation de la mémoire parcourue}

Dans le relèvement
\[
 201101|121221|102011|012222|102011
\]
les troisième et cinquième blocs coïncident :
\[
 B_3=B_5=102011.
\]
Leurs préétats modulo \(27\) sont cependant
\[
 p_3=(25,11,7,17,8,16),
 \qquad
 p_5=(7,8,4,23,26,7),
\]
et les quotients ultérieurs :
\[
 q_3=(6,13,9,2,13,1),
 \qquad
 q_5=(15,0,3,19,23,25).
\]
On vérifie \(p_3\ne p_5\) et \(q_3\ne q_5\).

\begin{figure}[htbp]
\centering
\begin{tikzpicture}[
 node distance=11mm and 7mm,
 every node/.style={font=\small},
 block/.style={draw,rounded corners,minimum width=17mm,minimum height=8mm},
 state/.style={draw,dashed,rounded corners,align=center,inner sep=3pt},
 >=stealth
]
\node[block] (b1) {\(201101\)};
\node[block,right=of b1] (b2) {\(121221\)};
\node[block,right=of b2] (b3) {\(102011\)};
\node[block,right=of b3] (b4) {\(012222\)};
\node[block,right=of b4] (b5) {\(102011\)};
\draw[->] (b1)--(b2);
\draw[->] (b2)--(b3);
\draw[->] (b3)--(b4);
\draw[->] (b4)--(b5);
\node[state,below=of b3] (p3)
 {\(p_3=(25,11,7,17,8,16)\)\\
  \(q_3=(6,13,9,2,13,1)\)};
\node[state,below=of b5] (p5)
 {\(p_5=(7,8,4,23,26,7)\)\\
  \(q_5=(15,0,3,19,23,25)\)};
\draw[->] (p3)--(b3);
\draw[->] (p5)--(b5);
\draw[<->,thick] (b3) to[bend left=28]
 node[above]{même mot visible} (b5);
\end{tikzpicture}
\caption{Retour visible de \(102011\) avec renouvellement du préétat et du
quotient.}
\label{gfp:b:fig:ruptura-memoria-e}
\end{figure}

\begin{theorem}[Retour visible sans périodicité de l’état]
\label{gfp:b:thm:retorno-visible-no-periodico-e}
\(B_3=B_5\) est un retour de la projection visible, non une orbite périodique
de l’état enrichi.
\end{theorem}

\begin{proof}
Le retour de l’état complet exigerait l’égalité de toutes ses coordonnées.
Les inégalités du préétat et du quotient l’excluent. De même,
\(1828|1828\) ne commence pas une période décimale : après le deuxième bloc apparaît
\(4\), tandis qu’une continuation périodique exigerait \(1\).
\end{proof}

\section{Certification factorielle ultérieure de la section exponentielle coinductive}

Soit \(\mathcal A\) une algèbre commutative unitaire sur \(\mathbb Q\), et
soit \(P_t\in\mathcal A[[t]]\) le propagateur formel. La concaténation
temporelle impose
\[
 P_{s+t}=P_sP_t,
 \qquad
 P_0=1.
\]
Écrivons
\[
 P_t=\sum_{n=0}^{\infty}a_nt^n.
\]
L’unité fixe \(a_0=1\), et le transport élémentaire normalisé fixe
\(a_1=1\).

\begin{theorem}[Récurrence factorielle]
\label{gfp:b:thm:recurrencia-factorial-e}
\[
 (n+1)a_{n+1}=a_n
 \qquad(n\geq0),
\]
et, par conséquent,
\[
 a_n=\frac1{n!},
 \qquad
 P_t=\sum_{n=0}^{\infty}\frac{t^n}{n!}.
\]
En \(t=1\),
\[
 \boxed{P_1=e.}
\]
\end{theorem}

\begin{proof}
Le coefficient de \(s^nt\) dans \(P_{s+t}\) est
\[
 \binom{n+1}{n}a_{n+1}=(n+1)a_{n+1}.
\]
Dans \(P_sP_t\), il est \(a_na_1=a_n\). L’égalité des séries produit la
récurrence. Une récurrence à partir de \(a_0=1\) donne \(a_n=1/n!\).
\end{proof}

\begin{corollary}[Équation différentielle]
\[
 \frac{d}{dt}P_t=P_t,
 \qquad P_0=1.
\]
\end{corollary}

\begin{proof}
\[
 P_t'=\sum_{n\geq0}(n+1)a_{n+1}t^n
 =\sum_{n\geq0}a_nt^n=P_t.
\]
\end{proof}

\section{Emboîtement cofinal des encadrements rationnels et convergence vers la valeur publiée}

Soit
\[
 S_N=\sum_{n=0}^{N}\frac1{n!}.
\]

\begin{theorem}[Borne factorielle explicite]
\label{gfp:b:thm:cota-factorial-e}
Pour \(N\geq1\),
\begin{equation}
 \boxed{
 S_N<e<
 S_N+\frac{N+2}{(N+1)(N+1)!}.}
 \label{gfp:b:eq:horquilla-factorial-e}
\end{equation}
\end{theorem}

\begin{proof}
Le reste est
\[
 R_N=\sum_{k=0}^{\infty}\frac1{(N+1+k)!}>0.
\]
Para \(k\geq0\),
\[
 (N+1+k)!
 \geq(N+1)!(N+2)^k.
\]
Par conséquent,
\[
 R_N
 \leq\frac1{(N+1)!}
 \sum_{k=0}^{\infty}\frac1{(N+2)^k}
 =\frac{N+2}{(N+1)(N+1)!}.
\]
L’inégalité est stricte à partir de \(k=2\).
\end{proof}

On définit
\[
 J_N^{(e)}=
 \left[
 S_N,\,
 S_N+\frac{N+2}{(N+1)(N+1)!}
 \right]
\]
et, pour la partie fractionnaire,
\[
 \widetilde J_N^{(e)}=J_N^{(e)}-2.
\]
La largeur converge vers zéro. Pour chaque \(M\), il existe \(N\) tel que
\[
 \operatorname{diam}J_N^{(e)}<10^{-M}.
\]
Si les extrémités appartiennent à une même cellule décimale, les \(M\) premiers
chiffres sont certifiés par arithmétique rationnelle.

\Needspace{30\baselineskip}
\section{Prolongement coinductif de la section exponentielle par les sous-cylindres ternaires du Topological Phase Kernel}

Au niveau \(K\), TPK et la relation \(\operatorname{Ext}\) ont déjà engendré
un unique sous-cylindre compatible. On prend le plus petit \(N=N_e(K)\) pour lequel
l’encadrement factoriel est contenu dans ce sous-cylindre. \(N_e(K)\) est un
ordre de certification analytique, non un indice générateur ni une règle de
sélection de la fibre.

Les cylindres TPK sont compatibles par troncature indépendamment du
semi-groupe. Le théorème du semi-groupe identifie ensuite leur évaluation
archimédienne à \(e-2\) ; la partie entière retrouve \(e\). Si l’orbite,
\(\operatorname{Ext}\) ou la troncature échouent, la génération échoue ; si la
récurrence ou la borne factorielle échouent, c’est ce certificat analytique qui échoue, non la
généalogie APP--TRIT--TPK.

\section{Certificat interne d’unicité orbitale, de récurrence factorielle, de convergence et de prolongement cylindrique}

\begin{remark}[Certificat et contrôle de validité : réfutation du canal exponentiel]
La construction est réfutée si :
\begin{enumerate}
 \item une autre orbite satisfait le prédicat fini complet ;
 \item la signature \(7[1828\,1828]459\) n’est pas positionnellement unique ;
 \item les préétats ou les quotients des deux blocs \(102011\) coïncident ;
 \item le semi-groupe ne force pas \((n+1)a_{n+1}=a_n\) ;
 \item la borne \eqref{gfp:b:eq:horquilla-factorial-e} échoue ;
 \item un niveau ne localise pas de manière unitaire une extension compatible ;
 \item un développement tabulé intervient avant la reconnaissance.
\end{enumerate}
\end{remark}
